\section{Metodologia}
\label{sec:metodologia}

\subsection{Organização do código}

O código foi dividido em três módulos principais de DSU em C++17, sem bibliotecas prontas para a lógica de conjuntos disjuntos: \texttt{DsuNaive}, \texttt{DsuUnionRank} e \texttt{DsuTarjan}, organizados em arquivos separados nos diretórios \texttt{include/} e \texttt{src/}. Além deles, o projeto mantém componentes compartilhados com funções bem definidas: \texttt{operacoes} gera as sequências reproduzíveis de teste, \texttt{estatisticas\_acesso} concentra a instrumentação de leituras e escritas, e \texttt{experimentos} executa o mesmo protocolo de medição para todas as variantes. Essa separação foi adotada para que uma mudança em uma implementação de DSU não exigisse alterar o gerador das operações nem a rotina de coleta das métricas, preservando a comparabilidade entre os experimentos.

\subsection{Instrumentação}

As leituras e escritas nos vetores internos da estrutura são contabilizadas por funções auxiliares que atualizam contadores por instância. Na variante ingênua, apenas o vetor \texttt{parent} entra nessa contagem; nas demais, contam-se acessos a \texttt{parent} e a \texttt{rank}. As inicializações realizadas antes da medição ficam fora desses totais, de modo que os valores reportados reflitam apenas o custo da sequência de operações \texttt{Find} e \texttt{Union}. O tempo de execução, por sua vez, é medido com \texttt{std::chrono::steady\_clock} ao redor do laço que reaplica a mesma sequência para cada variante.

\subsection{Cenários experimentais}

Para cada valor de $n$ na lista $\{500, 1000, 2000, 4000, 8000, 16000, 32000, 64000\}$, é gerada uma sequência com $m=20n$ operações por meio de \texttt{std::mt19937}. A fração de operações \texttt{Union} foi fixada em $0{,}25$, enquanto as demais correspondem a chamadas de \texttt{Find} em um único índice. Os índices são sorteados uniformemente em $\{0,\ldots,n-1\}$, e a semente usada para o $i$-ésimo valor de $n$ é $42 + 7919i$ (com $i=0$ para $n=500$), exatamente como no modo \texttt{--all} de \texttt{src/experimentos.cpp}. A mesma sequência é então reaplicada, na mesma ordem, às três variantes. Antes das medições, executa-se um aquecimento (\emph{warm-up}) com \texttt{DsuTarjan} sobre essa sequência, sem registrar seu tempo no CSV.

\subsection{Ambiente e compilação}

Os experimentos foram executados em um computador com processador Apple M1, 8\,GB de memória RAM e sistema operacional macOS 26.3 (build 25D125). A compilação foi feita com \texttt{g++}, correspondente ao Apple clang 17.0.0 (clang-1700.0.13.5), usando as flags \texttt{-std=c++17 -O3 -DNDEBUG}. O executável \texttt{bin/benchmark} grava os dados brutos em \texttt{results/benchmark.csv}. Em seguida, o script \texttt{scripts/export\_latex\_plot\_data.py} converte esses dados para \texttt{artigo/data/tempo\_por\_n.dat}, \texttt{artigo/data/acessos\_por\_n.dat} e \texttt{artigo/data/tabela\_resumo.inc.tex}, que são lidos diretamente pelo \LaTeX{} com \texttt{pgfplots}. As mesmas informações de ambiente ficam registradas em \texttt{docs/AMBIENTE.md} no repositório do projeto.
